\documentclass[11pt]{amsart}

\usepackage[margin=1in]{geometry}
\usepackage{amsmath,amssymb,amsthm,mathtools}
\usepackage{hyperref}
\usepackage[shortlabels]{enumitem}
\newtheorem{theorem}{Theorem}[section]
\newtheorem{proposition}[theorem]{Proposition}
\newtheorem{lemma}[theorem]{Lemma}
\newtheorem{corollary}[theorem]{Corollary}

\theoremstyle{definition}
\newtheorem{definition}[theorem]{Definition}

\theoremstyle{remark}
\newtheorem{remark}[theorem]{Remark}

\title{Local Stable Phase Retrieval in Fock Space}
\author{Jaume de Dios Pont}
\author{Mitchell Taylor}

\begin{document}
\begin{abstract}
We prove that the Bargmann--Fock space admits local stable phase retrieval at
the constant function, i.e. that there exists a constant \(\widetilde M\)
such that for all \(F\in \mathcal{F}^2(\mathbb{C})\) we have
\[
\inf_{\lvert \lambda \rvert=1}\lVert F-\lambda\mathbf 1 \rVert_{\mathcal{F}}
\le
\widetilde M\,\bigl\|\,\lvert F \rvert-1\,\bigr\|_{\mathcal{F}}.
\]
The proof proceeds in three stages. First, a general local reduction in
reproducing kernel Hilbert spaces shows that it suffices to prove the
statement for those \(F\) satisfying
\(\inf_{\lvert \lambda \rvert=1}\lVert F-\lambda\mathbf 1 \rVert_{\mathcal{F}}\ll1\). Second, a
general orthogonal reduction in Hilbert space shows that local stability near
any unit vector \(f_0\) follows from an ``orthogonal coercivity'' estimate for
perturbations perpendicular to \(f_0\). Finally, we establish this orthogonal
estimate for the Fock space at \(f_0=\mathbf 1\) with an explicit constant
\(C=4600\), by reducing the Gaussian integral to a family of circle problems
via polar coordinates, then controlling each circle uniformly using a
combination of low-frequency and high-frequency arguments. The entire proof
has been formalised in Lean~4.
\end{abstract}
\maketitle

\section{Introduction}\label{sec:intro}

\subsection{Phase retrieval and stability}

Phase retrieval asks whether a signal can be recovered from the magnitudes of
its measurements alone.  The problem arises throughout science and
engineering---in crystallography, coherent diffraction imaging, and audio
processing---whenever detectors record intensities but not phases.
Mathematically, one seeks to invert (up to unavoidable phase ambiguities) the
map $f \mapsto \lvert f \rvert$ or, more generally, $f \mapsto \lvert Tf \rvert$ for a
suitable measurement operator~$T$.

In finite dimensions the theory is by now extensive; see the
survey~\cite{grohs2020survey}.  In infinite dimensions the situation is
considerably more delicate.  A fundamental negative result of Alaifari and
Grohs~\cite{alaifari2017continuous} shows that phase retrieval for continuous
frames on infinite-dimensional Hilbert spaces is always unstable: the local
Lipschitz constant of the reconstruction map is necessarily unbounded.  For
the Gabor transform (the short-time Fourier transform with a Gaussian
window), Grohs and Rathmair~\cite{grohs2020survey} discovered a deep
connection between the stability constant and spectral geometry, showing that
the reciprocal of the Cheeger constant of the time--frequency landscape
$\lvert V_\varphi f \rvert$ controls the local stability of reconstruction.  More
recently, Alaifari, Pineau, Taylor, and
Wellershoff~\cite{alaifari2021gabor} proved that for any reasonable
window function and any weighted Sobolev norm up to the sharp regularity
threshold, the local stability constant is infinite on a dense set---confirming
that instability is pervasive for short-time Fourier transform phase retrieval.

In a complementary direction, Calderbank, Daubechies, Freeman, and
Freeman~\cite{calderbank2022stable} showed that stable phase retrieval
\emph{is} possible in infinite dimensions for operators that do not arise from
frames, and Christ, Pineau, and Taylor~\cite{christ2022examples} constructed
explicit infinite-dimensional subspaces of $L^2(\mathbb{R})$ admitting H\"older-stable
phase retrieval.  The function-space framework of Freeman, Oikhberg, Pineau,
and Taylor~\cite{FOTP} systematised these developments: given a closed
subspace $E$ of a complex Banach lattice~$X$, one says that $E$ does
\emph{stable phase retrieval} with constant~$C$ if
\[
\inf_{\lvert w \rvert=1}\lVert f - wg \rVert_X
\;\le\;
C\,\bigl\|\,\lvert f \rvert - \lvert g \rvert\,\bigr\|_X
\qquad
\text{for all } f,g \in E.
\]
Here the infimum over unit complex numbers~$w$ accounts for the inherent
phase ambiguity: the map $f \mapsto \lvert f \rvert$ cannot distinguish $f$ from
$wf$.  The results of~\cite{FOTP} show that many order-continuous Banach
lattices admit no infinite-dimensional subspace with this property, but leave
open whether specific structured spaces---particularly spaces of entire
functions---might enjoy stable phase retrieval, at least locally.

\subsection{The Bargmann--Fock space}

For any measurable $f \colon \mathbb{C} \to \mathbb{C}$, define the
\emph{Fock norm}
\[
\lVert f \rVert_{\mathcal{F}}^2
\;:=\;
\frac{1}{\pi}\int_{\mathbb{C}} \lvert f(z) \rvert^2 \, \mathrm{e}^{-\lvert z \rvert^2}\,\mathrm{d} m(z).
\]
The \emph{Bargmann--Fock space}~\cite{bargmann1961hilbert}
$\mathcal{F}^2(\mathbb{C})$ is the closed subspace of entire functions with
$\lVert f \rVert_\mathcal{F} < \infty$.
The monomials $e_n(z) = z^n/\sqrt{n!}$ form an orthonormal basis, so every
$f(z) = \sum_{n \ge 0} a_n z^n$ satisfies
$\lVert f \rVert_{\mathcal{F}}^2 = \sum_{n \ge 0} \lvert a_n \rvert^2 n!$.
With this convention the Fock norm is defined for arbitrary measurable
functions, not only holomorphic ones; in particular
$\lVert \,\lvert f \rvert\, \rVert_\mathcal{F} = \lVert f \rVert_\mathcal{F}$.

The Fock space occupies a special place in the phase retrieval landscape.  Its
elements are entire functions, hence determined by their values on any open
set, and the Gaussian weight provides strong localisation of
high-degree monomials.  These structural features make it a natural candidate
for stable phase retrieval results that fail in more general $L^2$ settings.

\subsection{The main result}

Given a unit-norm element $f_0 \in \mathcal{F}^2(\mathbb{C})$ and a small perturbation $p$,
we cannot hope to recover the perturbation~$p$ from the magnitude
$\lvert f_0 + p \rvert$ alone: for any unit complex number $w$, we have
$\lvert w(f_0 + p) \rvert = \lvert f_0 + p \rvert$, so the magnitude measurement cannot
distinguish $f_0 + p$ from its phase rotation $w(f_0 + p)$.  The correct
question is whether, after optimising over the phase~$w$, the Fock norm of
$w(f_0+p) - f_0$ is controlled by the ``magnitude residual''
$\bigl\|\, \lvert f_0 + p \rvert - \lvert f_0 \rvert\,\bigr\|_\mathcal{F}$.

\begin{theorem}[Local stable phase retrieval in Fock space]\label{thm:local}
\par\noindent\texttt{FockSPR.LocalFockSPR\_exists\_phase}\par\smallskip
There exist absolute constants $\delta > 0$ and $\widetilde{M} > 0$ such
that the following holds.  Let $p$ be a polynomial with
$\lVert p \rVert_\mathcal{F} \le \delta$.  Then there exists $w \in \mathbb{C}$ with $\lvert w \rvert = 1$
such that
\[
\lVert w(1+p) - 1 \rVert_\mathcal{F}
\;\le\;
\widetilde{M}\,
\bigl\|\, \lvert 1+p \rvert - 1 \,\bigr\|_{\mathcal{F}}.
\]
One may take $\delta = 1/4601$ and $\widetilde{M} = 23003$.
\end{theorem}

\begin{remark}\label{rem:wlog}
We state the result for $f_0 = \mathbf{1}$ (the constant function~$1$).
This is not without loss of generality: the Fock space does not possess an
isometry mapping an arbitrary unit-norm element to the constant function.
However, the constant function is a natural base point because the orthogonal
complement $\mathbf{1}^\perp$ consists precisely of the functions vanishing at
the origin, and the monomial structure of the Fock space makes this condition
especially tractable.  The condition that $p$ is a polynomial can be
relaxed to $1 + p \in \mathcal{F}^2(\mathbb{C})$ by approximation.
\end{remark}

Theorem~\ref{thm:local} is deduced from the following ``orthogonal
coercivity'' estimate, which is the technical heart of the paper:

\begin{theorem}[Orthogonal coercivity]\label{thm:orthog}
\par\noindent\texttt{LocalFockSPR}\par\smallskip
Let $U(z) = \sum_{n=1}^D a_n z^n$ be a polynomial with $U(0) = 0$.  Then
\[
\lVert U \rVert_\mathcal{F}^2
\;\le\;
4600^2 \cdot
\lVert \rho \circ U \rVert_\mathcal{F}^2,
\]
where $\rho(w) = \bigl\lvert\, \lvert 1+w \rvert - 1 \,\bigr\rvert$.
\end{theorem}

The condition $U(0) = 0$ means that $U$ is orthogonal to the constant
function $\mathbf{1}$ in $\mathcal{F}^2(\mathbb{C})$.  The passage from
Theorem~\ref{thm:orthog} to Theorem~\ref{thm:local} is a general
Hilbert-space argument, carried out in Section~\ref{sec:reduction}.  The
proof of Theorem~\ref{thm:orthog}---the content of
Sections~\ref{sec:circle}--\ref{sec:assembly}---is where the analysis
specific to Fock space takes place.  When combined with the local reduction
of Section~\ref{sec:local-reduction}, we arrive at the following main
theorem.

\begin{theorem}[Main theorem]\label{thm:main}
There exists an absolute constant \(M>0\) such that for all
\(F\in\mathcal{F}^2(\mathbb{C})\) there exists \(w\in\mathbb{C}\) with \(\lvert w \rvert=1\) such that
\[
\lVert wF-\mathbf 1 \rVert_{\mathcal{F}}
\le
M\,\bigl\|\,\lvert F \rvert-1\,\bigr\|_{\mathcal{F}}.
\]
\end{theorem}

\subsection{Strategy of proof}\label{sec:strategy}

\noindent
\textbf{The local-local to local reduction
(Section~\ref{sec:local-reduction}).}\;
A general argument in reproducing kernel Hilbert spaces shows that if one
can prove phase retrieval at a function \(f_0\) and establish local stability
with respect to very small perturbations of \(f_0\), then one can prove local
stability for general perturbations of \(f_0\).

\medskip
\noindent
\textbf{Orthogonal reduction (Section~\ref{sec:reduction}).}\;
A general argument in Hilbert space shows that if orthogonal perturbations
of a unit vector $f_0$ satisfy a coercivity estimate with constant~$M$,
then all sufficiently small perturbations (after optimising over a global
phase) satisfy a coercivity estimate with constant $\widetilde{M} = O(M)$.
The key ingredients are the orthogonal decomposition $h = af_0 + g$ with
$g \perp f_0$, control of the scalar coefficient $a$ via a Cauchy--Schwarz
argument, and a smallness absorption.

\medskip
\noindent
\textbf{Polar coordinates.}\;
The Gaussian integral decomposes as
\begin{equation}\label{eq:polar}
\frac{1}{\pi}\int_{\mathbb{C}} F(z)\, \mathrm{e}^{-\lvert z \rvert^2}\,\mathrm{d} m(z)
\;=\;
2\int_0^\infty r\,\mathrm{e}^{-r^2}
\biggl(\int_{S^1} F(r\mathrm{e}^{\mathrm{i} t})\,\mathrm{d}\sigma(t)\biggr)\mathrm{d} r,
\end{equation}
where $\mathrm{d}\sigma = \mathrm{d} t/(2\pi)$ is the probability Haar measure on the
circle.  Both norms in Theorem~\ref{thm:orthog} decompose in this way, so it
suffices to compare the inner circle integrals for each radius~$r$.

\medskip
\noindent
\textbf{Circle estimates (Section~\ref{sec:circle}).}\;
On a circle of radius~$r$, the function $t \mapsto U(r\mathrm{e}^{\mathrm{i} t})$ is a
trigonometric polynomial with strictly positive frequencies (since
$U(0) = 0$).  We prove two estimates relating the $L^2$ norm of such a
polynomial~$P$ to that of $\rho \circ P$:
\begin{enumerate}[(i)]
\item A \emph{local circle estimate} that works for any finite set of
  positive frequencies, at the cost of a constant depending on how many
  frequencies are present.

\item A \emph{high-frequency estimate} that applies when the lowest
  frequency is much larger than the bandwidth, giving a universal constant
  independent of the number of frequencies.
\end{enumerate}

\medskip
\noindent
\textbf{Frequency localisation (Section~\ref{sec:blocks}).}\;
The polynomial $U$ may have up to $D$ frequencies, but the monomial~$z^n$
concentrates its Gaussian mass near the annulus $\lvert z \rvert \approx \sqrt{n}$.
We exploit this by partitioning the frequencies into blocks, each localised
to a specific annulus.  For each annulus, only a bounded window of nearby
blocks contributes significantly; the energy leaking from distant blocks
decays as a Gaussian in the block distance.

\medskip
\noindent
\textbf{Assembly (Section~\ref{sec:assembly}).}\;
On each annulus, the nearby blocks form a trigonometric polynomial to which
one of the two circle estimates applies.  This gives a uniform bound on the
$\rho$-norm of the ``local'' polynomial.  We then use the $1$-Lipschitz
property of $\rho$ to pass from the local polynomial to the full $U$,
integrate over all annuli, and absorb the exponentially small leakage.

\subsection{Notation}\label{sec:notation}

All circle integrals use the probability measure
$\mathrm{d}\sigma = \mathrm{d} t/(2\pi)$, so that Parseval reads
$\lVert f \rVert_{L^2}^2 = \sum_n \lvert \hat{f}(n) \rvert^2$.
The Lean theorem names recorded in blue monospace throughout the
paper refer to the accompanying Lean~4 formalisation, which uses the Mathlib
library~\cite{mathlib}.  The only things that have not been Lean verified are
the local-local to local reduction and the passage from perturbations which
are polynomials to general perturbations in the Fock space. Full details are
included.

\section{The local reduction}\label{sec:local-reduction}

In this section we show that local stable phase retrieval is implied by an
even more local variant.

\begin{proposition}[Local reduction]\label{prop:local-reduction}
Let \(H\) be a reproducing kernel Hilbert space of real- or complex-valued
functions isomorphically embedded in \(L^2(\Omega,\mu)\). Fix \(F\in H\) and
\(\epsilon>0\). Suppose that there exists a constant \(C>0\) such that for
all \(G\in H\) satisfying
\[
\inf_{\lvert \lambda \rvert=1}\lVert F-\lambda G \rVert_{L^2(\Omega)}<\epsilon
\]
we have
\[
\inf_{\lvert \lambda \rvert=1}\lVert F-\lambda G \rVert_{L^2(\Omega)}
\le
C\lVert \lvert F \rvert-\lvert G \rvert \rVert_{L^2(\Omega)}.
\]
Suppose also that phase retrieval holds at \(F\), i.e. if \(G\in H\) and
\(\lvert F \rvert=\lvert G \rvert\), then there exists \(\lambda\in\mathbb{C}\) with
\(\lvert \lambda \rvert=1\) and \(F=\lambda G\). Then there exists a constant
\(C'>0\) such that for all \(G\in H\),
\[
\inf_{\lvert \lambda \rvert=1}\lVert F-\lambda G \rVert_{L^2(\Omega)}
\le
C'\lVert \lvert F \rvert-\lvert G \rvert \rVert_{L^2(\Omega)}.
\]
\end{proposition}

\begin{proof}
We need to verify that there exists a constant \(C'\) such that for all
\(G\in H\),
\[
\inf_{\lvert \lambda \rvert=1}\lVert F-\lambda G \rVert_{L^2(\Omega)}
\le
C'\lVert \lvert F \rvert-\lvert G \rvert \rVert_{L^2(\Omega)}.
\]

First suppose that \(\lVert G \rVert_{L^2(\Omega)}\ge4\lVert F \rVert_{L^2(\Omega)}\).
Then
\[
\lVert \lvert F \rvert-\lvert G \rvert \rVert_{L^2(\Omega)}
\ge
\lVert G \rVert_{L^2(\Omega)}-\lVert F \rVert_{L^2(\Omega)}
\ge
\frac34\lVert G \rVert_{L^2(\Omega)},
\]
while
\[
\inf_{\lvert \lambda \rvert=1}\lVert F-\lambda G \rVert_{L^2(\Omega)}
\le
\lVert F \rVert_{L^2(\Omega)}+\lVert G \rVert_{L^2(\Omega)}
\le
\frac54\lVert G \rVert_{L^2(\Omega)}.
\]
Thus the desired inequality holds in this case with any \(C'\ge5/3\).
Therefore we may restrict attention to \(G\) with
\(\lVert G \rVert_{L^2(\Omega)}<4\lVert F \rVert_{L^2(\Omega)}\).

Now suppose that the desired inequality is true for all \(G\) satisfying
\[
\inf_{\lvert \lambda \rvert=1}\lVert F-\lambda G \rVert_{L^2(\Omega)}<\epsilon,
\]
but is not true in general. Then for every \(n\) we can find \(G_n\in H\)
with
\[
\inf_{\lvert \lambda \rvert=1}\lVert F-\lambda G_n \rVert_{L^2(\Omega)}\ge\epsilon,
\]
\[
\inf_{\lvert \lambda \rvert=1}\lVert F-\lambda G_n \rVert_{L^2(\Omega)}
>
n\lVert \lvert F \rvert-\lvert G_n \rvert \rVert_{L^2(\Omega)},
\]
and also \(\lVert G_n \rVert_{L^2(\Omega)}<4\lVert F \rVert_{L^2(\Omega)}\). This latter
inequality guarantees that
\(\inf_{\lvert \lambda \rvert=1}\lVert F-\lambda G_n \rVert_{L^2(\Omega)}\) is uniformly
bounded, and hence
\[
\lVert \lvert F \rvert-\lvert G_n \rvert \rVert_{L^2(\Omega)}\to0.
\]
In particular, after passing to a subsequence, we may assume that
\[
\lvert G_{n_k}(x) \rvert\to\lvert F(x) \rvert
\qquad\text{for almost every }x\in\Omega.
\]

Since \(\lVert G_{n_k} \rVert_{L^2(\Omega)}\) is bounded independently of \(k\), and
the embedding of \(H\) into \(L^2(\Omega,\mu)\) is isomorphic, the sequence
\((G_{n_k})\) is bounded in \(H\). Passing to a further subsequence if
necessary, we may assume that \(G_{n_{k_\ell}}\rightharpoonup G\) weakly in
\(H\). Since point evaluations are bounded on \(H\), we have
\[
G_{n_{k_\ell}}(x)\to G(x)
\]
for every \(x\in\Omega\). Hence \(\lvert G_{n_{k_\ell}}(x) \rvert\to\lvert G(x) \rvert\) for
every \(x\in\Omega\). This forces \(\lvert G \rvert=\lvert F \rvert\) almost everywhere. By the
phase-retrieval hypothesis, \(G=\lambda F\) for some unimodular \(\lambda\).

Collecting what we know, \(G_{n_{k_\ell}}\) converges weakly in \(H\) to
\(\lambda F\). Moreover, from
\(\lVert \lvert F \rvert-\lvert G_n \rvert \rVert_{L^2(\Omega)}\to0\), we see that
\[
\lVert G_n \rVert_{L^2(\Omega)}=\lVert \lvert G_n \rvert \rVert_{L^2(\Omega)}
\to
\lVert \lvert F \rvert \rVert_{L^2(\Omega)}
=\lVert \lambda F \rVert_{L^2(\Omega)}.
\]
Because the embedding is isomorphic, weak convergence in \(H\) together with
convergence of the \(L^2\)-norms to the norm of the weak limit implies strong
convergence in \(L^2\). Thus \(G_{n_{k_\ell}}\to\lambda F\) in
\(L^2(\Omega)\). For all sufficiently large \(\ell\), this contradicts
\[
\inf_{\lvert \lambda \rvert=1}
\lVert F-\lambda G_{n_{k_\ell}} \rVert_{L^2(\Omega)}
\ge \epsilon.
\]
\end{proof}

\section{The orthogonal reduction}\label{sec:reduction}

In this section we work in the generality of a complex Hilbert space $H$
continuously embedded in $L^2(\mu)$ for some measure space $(\Omega, \mu)$.
We show that \emph{orthogonal coercivity}---a stability estimate for
perturbations perpendicular to a fixed unit vector---implies \emph{local
stability} for all sufficiently small perturbations, after optimising over a
global phase.  The argument is purely Hilbert-space-theoretic and uses no
structure specific to the Fock space.

\begin{theorem}[Orthogonal reduction]\label{thm:reduction}
\par\noindent\texttt{FockSPR.LocalFockSPR\_constants}\par\smallskip
Let $H$ be a complex Hilbert space with norm $\lVert \cdot \rVert$, continuously
embedded in $L^2(\mu)$.  Let $f_0 \in H$ with $\lVert f_0 \rVert = 1$, and
suppose there exists $M > 0$ such that
\begin{equation}\label{eq:orthog-hyp}
\lVert g \rVert
\;\le\;
M\, \bigl\|\, \lvert f_0 + g \rvert - \lvert f_0 \rvert \,\bigr\|_{L^2}
\qquad
\text{for all } g \perp f_0.
\end{equation}
Then there exist $\delta > 0$ and $\widetilde{M} > 0$, depending only
on~$M$, such that:
for every $h \in H$ with $\lVert h \rVert \le \delta$ and
$\langle h,f_0\rangle \in \mathbb{R}$, we have
\[
\lVert h \rVert
\;\le\;
\widetilde{M}\, \bigl\|\, \lvert f_0 + h \rvert - \lvert f_0 \rvert \,\bigr\|_{L^2}.
\]
One may take $\delta = 1/(M+1)$ and $\widetilde{M} = 5M + 3$.
\end{theorem}

\begin{remark}
The condition $\langle h,f_0\rangle \in \mathbb{R}$ is a phase normalisation: it says that
$h$ has been rotated so that its component along $f_0$ is real.  In the
application to Theorem~\ref{thm:local}, this is achieved by choosing the
optimal unit phase $w$ so that $\langle w(f_0 + p) - f_0,f_0\rangle$ is real and
nonneg\-ative, which amounts to aligning the phase of the inner product
$\langle p,f_0\rangle$.
\end{remark}

\begin{proof}
Write $m = \bigl\|\, \lvert f_0 + h \rvert - \lvert f_0 \rvert \,\bigr\|_{L^2}$ for the
magnitude residual.

\medskip
\noindent
\emph{Step 1: Orthogonal decomposition.}\;
Set $a = \operatorname{Re}\langle h,f_0\rangle = \langle h,f_0\rangle$ (a real number by hypothesis) and
$g = h - a f_0$.  Then $g \perp f_0$, $h = af_0 + g$, and
$\lVert h \rVert^2 = a^2 + \lVert g \rVert^2$, so
\begin{equation}\label{eq:triangle}
\lVert h \rVert \;\le\; \lvert a \rvert + \lVert g \rVert.
\end{equation}

\medskip
\noindent
\emph{Step 2: Control of the scalar coefficient.}\;
The pointwise identity
$\lvert f_0 + h \rvert^2 - \lvert f_0 \rvert^2
= 2\,\operatorname{Re}(h\,\overline{f_0}) + \lvert h \rvert^2$
integrates to
\[
\int_\Omega \bigl(\lvert f_0 + h \rvert^2 - \lvert f_0 \rvert^2\bigr)\,\mathrm{d}\mu
= 2a + \lVert h \rVert^2.
\]
On the other hand, the factorisation
$\lvert f_0 + h \rvert^2 - \lvert f_0 \rvert^2
= \bigl(\lvert f_0 + h \rvert - \lvert f_0 \rvert\bigr)
  \bigl(\lvert f_0 + h \rvert + \lvert f_0 \rvert\bigr)$
and Cauchy--Schwarz give
\[
\lvert 2a + \lVert h \rVert^2 \rvert
\;\le\;
m \cdot \bigl\|\, \lvert f_0 + h \rvert + \lvert f_0 \rvert \,\bigr\|_{L^2}
\;\le\;
m\,(2 + \lVert h \rVert).
\]
Hence
\begin{equation}\label{eq:a-bound}
\lvert a \rvert
\;\le\;
\frac{2 + \lVert h \rVert}{2}\, m
+ \frac{\lVert h \rVert^2}{2}.
\end{equation}

\medskip
\noindent
\emph{Step 3: Control of the orthogonal component.}\;
The reverse triangle inequality
$\bigl\lvert \lvert f_0 + g \rvert - \lvert f_0 + h \rvert \bigr\rvert
\le \lvert g - h \rvert = \lvert a \rvert\,\lvert f_0 \rvert$
gives
$\bigl\|\, \lvert f_0 + g \rvert - \lvert f_0 \rvert \,\bigr\|_{L^2}
\le \lvert a \rvert + m$.
The orthogonal coercivity hypothesis~\eqref{eq:orthog-hyp} then yields
\begin{equation}\label{eq:g-bound}
\lVert g \rVert \;\le\; M\,(\lvert a \rvert + m).
\end{equation}

\medskip
\noindent
\emph{Step 4: Combine.}\;
Substituting~\eqref{eq:a-bound} and~\eqref{eq:g-bound}
into~\eqref{eq:triangle}:
\[
\lVert h \rVert
\;\le\;
\lvert a \rvert + M(\lvert a \rvert + m)
= (M+1)\lvert a \rvert + Mm
\;\le\;
\underbrace{\Bigl(M + \frac{(M+1)(2+\lVert h \rVert)}{2}\Bigr)}_{=:\, C(\lVert h \rVert)}
m
\;+\; \frac{M+1}{2}\,\lVert h \rVert^2.
\]

\medskip
\noindent
\emph{Step 5: Smallness absorption.}\;
Since $\lVert h \rVert \le \delta$, we have $\lVert h \rVert^2 \le \delta\,\lVert h \rVert$, so
\[
\lVert h \rVert
\;\le\; C(\delta)\, m + \frac{(M+1)\delta}{2}\,\lVert h \rVert.
\]
Choosing $\delta = 1/(M+1)$ makes $(M+1)\delta/2 = 1/2$, giving
$\lVert h \rVert/2 \le C(\delta)\, m$ and therefore
\[
\lVert h \rVert \;\le\; 2\, C(\delta)\, m.
\]
A short computation shows
$C(\delta) = M + (M+1)(2 + \delta)/2 < (5M+3)/2$, so one may take
$\widetilde{M} = 5M + 3$.
\end{proof}

\begin{corollary}[Fock-space application]\label{cor:fock-local}
Taking $f_0 = \mathbf{1}$ (the constant function), $H = \mathcal{F}^2(\mathbb{C})$, and
$M = 4600$ (from Theorem~\ref{thm:orthog}), the orthogonal reduction gives
$\delta = 1/4601$ and $\widetilde{M} = 23003$.
This is Theorem~\ref{thm:local}.
\end{corollary}

\begin{proof}
The condition $g \perp \mathbf{1}$ in $\mathcal{F}^2(\mathbb{C})$ means precisely that
$g(0) = 0$, i.e.\ $g$ has no constant term.
Theorem~\ref{thm:orthog} provides the orthogonal coercivity
hypothesis~\eqref{eq:orthog-hyp} with $M = 4600$.
The optimal phase $w$ is chosen so that
$h = w(1 + p) - 1$ satisfies $\langle h,\mathbf{1}\rangle \in \mathbb{R}$ and
$\lVert h \rVert_\mathcal{F} \le \delta$; the existence of such a $w$ is straightforward.
Theorem~\ref{thm:reduction} then yields the conclusion of
Theorem~\ref{thm:local}.
\end{proof}

 \section{Circle estimates}\label{sec:circle}

In this section we prove two estimates controlling the $L^2$ norm of a
trigonometric polynomial $P$ on the circle in terms of the $L^2$ norm of
$\rho \circ P$.  Both exploit the fact that all frequencies of~$P$ are
strictly positive---a consequence of $U(0) = 0$---but they apply in
complementary regimes.  The \emph{local circle estimate}
(Theorem~\ref{thm:local-circle}) works for any frequency set and gives a
constant proportional to $\sqrt{L}$, where $L$ is the number of frequencies.
The \emph{high-frequency estimate} (Theorem~\ref{thm:high-freq}) applies
when the lowest frequency is much larger than the bandwidth, and gives an
absolute constant.  Together they will provide a uniform bound on every
annulus.

Throughout this section, $P$ denotes a trigonometric polynomial
$P(t) = \sum_{n \in E} b_n \mathrm{e}^{\mathrm{i} nt}$ with $E \subset \mathbb{Z}_{>0}$ finite,
and all norms are $L^2(S^1, \mathrm{d}\sigma)$.

Recall the function $\rho(w) = \bigl\lvert \lvert 1+w \rvert - 1 \bigr\rvert$,
which is $1$-Lipschitz by two applications of the reverse triangle inequality:
$\lvert \rho(w) - \rho(z) \rvert
\le \bigl\lvert \lvert 1+w \rvert - \lvert 1+z \rvert \bigr\rvert
\le \lvert w - z \rvert$.
In particular, $\rho(w) \le \lvert w \rvert$ (take $z = 0$) and
\begin{equation}\label{eq:rho-lip}
\rho(w) \;\ge\; \rho(z) - \lvert w - z \rvert.
\end{equation}
\par\noindent\texttt{FockSPR.rho\_lipschitz}\par\smallskip

\subsection{The local circle estimate}\label{sec:local}

A preliminary inequality will be useful both here and in the high-frequency
estimate.

\begin{lemma}[Safe-square inequality]\label{lem:safe-square}
\par\noindent\texttt{FockSPR.nonneg\_safe\_square}\par\smallskip
Let $a, b, c \ge 0$ with $a \ge b - c$.  Then
$a^2 \ge \tfrac{1}{2}\, b^2 - c^2$.
\end{lemma}

\begin{proof}
If $b \le c$ the right side is nonpositive.  If $b > c$ then
$a \ge b-c > 0$ gives $a^2 \ge (b-c)^2 \ge \tfrac{1}{2}b^2 - c^2$, where
the last step is the identity
$(b-c)^2 - \tfrac{1}{2}b^2 + c^2 = \tfrac{1}{2}(b-2c)^2 \ge 0$.
\end{proof}

\begin{theorem}[Local circle estimate]\label{thm:local-circle}
\par\noindent\texttt{FockSPR.local\_circle\_estimate}\par\smallskip
Let $E \subset \mathbb{Z}_{>0}$ with $\lvert E \rvert = L \ge 1$, let $b_n \in \mathbb{C}$ for
$n \in E$, and set $P(t) = \sum_{n \in E} b_n\,\mathrm{e}^{\mathrm{i} nt}$.  Then
\[
\lVert P \rVert_{L^2}^2
\;\le\;
144\, L \cdot \lVert \rho \circ P \rVert_{L^2}^2.
\]
\end{theorem}

\begin{proof}
Write $\alpha = \lVert P \rVert_{L^2}$ and $\beta = \lVert \rho \circ P \rVert_{L^2}$.
We show $\beta \ge \alpha/(12\sqrt{L})$, which gives the result after
squaring.  The threshold between the two regimes is
$\alpha_0 = 1/(4\sqrt{L})$.

Two structural properties of $P$ underpin both cases.  First,
\emph{Cauchy--Schwarz and Parseval} give the pointwise bound
$\lvert P(t) \rvert \le \sqrt{L}\, \alpha$ for all~$t$.
Second, since every $n \in E$ satisfies $n \ge 1$, the zeroth Fourier
coefficient of $P$ vanishes:
$\int_{S^1} P\,\mathrm{d}\sigma = 0$.
More importantly, $n + m \ge 2$ for all $n, m \in E$, so
$\int_{S^1} P^2\,\mathrm{d}\sigma = 0$ as well.  Writing $P = u + \mathrm{i} v$ with
$u, v$ real-valued, this identity splits as
\begin{equation}\label{eq:equal-mass}
\int_{S^1} u^2\,\mathrm{d}\sigma
= \int_{S^1} v^2\,\mathrm{d}\sigma
= \frac{\alpha^2}{2}
\qquad\text{and}\qquad
\int_{S^1} uv\,\mathrm{d}\sigma = 0.
\end{equation}

\medskip
\noindent
\textbf{Small amplitudes} ($\alpha \le \alpha_0$).
Then $\lVert P \rVert_\infty \le \sqrt{L}\,\alpha \le 1/4$.  For $\lvert w \rvert \le 1/4$
we have $\lvert 1+w \rvert \ge 3/4$ and $\operatorname{Re}(w) \ge -1/4$, so the denominator
$\lvert 1+w \rvert + 1 + \operatorname{Re}(w) \ge 3/2 > 0$.  The algebraic identity
\[
\lvert 1+w \rvert - 1 - \operatorname{Re}(w)
= \frac{\operatorname{Im}(w)^2}{\lvert 1+w \rvert + 1 + \operatorname{Re}(w)}
\]
shows that $\lvert 1+w \rvert - 1 = \operatorname{Re}(w) + \varepsilon$ with
$0 \le \varepsilon \le \operatorname{Im}(w)^2/(3/2) \le \lvert w \rvert^2$.  In particular
$\rho(w) = \bigl\lvert \operatorname{Re}(w) + \varepsilon\bigr\rvert \ge
\lvert \operatorname{Re}(w) \rvert - \varepsilon$, so by the safe-square inequality
(Lemma~\ref{lem:safe-square}),
\[
\rho(P)^2 \;\ge\; \tfrac{1}{2}\, u^2 - \lvert P \rvert^4.
\]
Integrating and using~\eqref{eq:equal-mass}:
$\beta^2 \ge \alpha^2/4 - \lVert P \rVert_\infty^2\, \alpha^2 \ge \alpha^2/4 - L\,\alpha^4$.
Since $L\alpha^2 \le 1/16$, we get $\beta^2 \ge 3\alpha^2/16$, so
$\beta \ge \sqrt{3}\,\alpha/4 \ge \alpha/(12\sqrt{L})$.

\medskip
\noindent
\textbf{Large amplitudes} ($\alpha > \alpha_0$).
The vanishing of $\hat P(0)$ gives
$\int_{S^1}\lvert 1+P \rvert^2\,\mathrm{d}\sigma = 1 + \alpha^2$, so
\[
\alpha^2
= \int_{S^1}\bigl(\lvert 1+P \rvert^2 - 1\bigr)\,\mathrm{d}\sigma
\le \int_{S^1} \rho(P)\,\bigl(\lvert 1+P \rvert + 1\bigr)\,\mathrm{d}\sigma,
\]
where we used $\bigl\lvert \lvert 1+P \rvert^2 - 1\bigr\rvert
= \bigl(\lvert 1+P \rvert + 1\bigr)\rho(P)$.  By Cauchy--Schwarz,
\[
\alpha^2 \;\le\; \beta \cdot \bigl(\sqrt{1+\alpha^2} + 1\bigr).
\]
If $\alpha \le 1$: $\beta \ge \alpha^2/3$ and $\alpha > \alpha_0 = 1/(4\sqrt{L})$ give
$\beta \ge \alpha/(12\sqrt{L})$.
If $\alpha > 1$: $\beta \ge \alpha/3 \ge \alpha/(12\sqrt{L})$.

\medskip
In both regimes $\beta \ge \alpha/(12\sqrt{L})$.
\end{proof}


\subsection{Rotational averaging}\label{sec:rot-avg}

The local circle estimate grows with the number of frequencies~$L$.  To
eliminate this dependence for high-frequency polynomials, we will need the
following fact: the $\rho$-image of a full circle already captures a
definite fraction of its radius.

\begin{proposition}[Rotational averaging]\label{prop:rot-avg}
\par\noindent\texttt{FockSPR.rotational\_averaging\_bound}\par\smallskip
For all $r \ge 0$,
\[
\int_{S^1} \rho(r\mathrm{e}^{\mathrm{i}\theta})^2\,\mathrm{d}\sigma(\theta)
\;\ge\; \frac{r^2}{8}.
\]
\end{proposition}

\begin{proof}
On the arc $\lvert \theta \rvert \le \pi/4$ (which has $\mathrm{d}\sigma$-measure $1/4$),
we claim $\rho(r\mathrm{e}^{\mathrm{i}\theta}) \ge r/\sqrt{2}$.  Indeed,
$\cos\theta \ge 1/\sqrt{2}$ gives
\[
\lvert 1 + r\mathrm{e}^{\mathrm{i}\theta} \rvert - 1
= \frac{r^2 + 2r\cos\theta}{\lvert 1+r\mathrm{e}^{\mathrm{i}\theta} \rvert + 1}
\ge \frac{r(r + \sqrt{2})}{r + 2}
\ge \frac{r}{\sqrt{2}},
\]
where the last inequality is $(\sqrt{2}-1)r \ge 0$.  Moreover the numerator
is positive, so $\lvert 1 + r\mathrm{e}^{\mathrm{i}\theta} \rvert > 1$ and $\rho$ equals
$\lvert 1+r\mathrm{e}^{\mathrm{i}\theta} \rvert - 1$ without absolute value.  Restricting the
integral to this arc:
$\int \rho^2\,\mathrm{d}\sigma \ge \frac{1}{4}(r/\sqrt{2})^2 = r^2/8$.
\end{proof}


\subsection{The high-frequency estimate}\label{sec:high-freq}

When the polynomial oscillates rapidly, the Lipschitz property
of~$\rho$ and the rotational averaging of
Proposition~\ref{prop:rot-avg} combine to give an absolute bound.  The idea
is to factor $P(t) = \mathrm{e}^{\mathrm{i} Nt}\, Q(t)$ where $Q$ is a slowly varying
envelope; on short intervals the leading phase $\mathrm{e}^{\mathrm{i} Nt}$ sweeps through
a full rotation while $Q$ barely changes, and the rotational averaging
ensures that $\rho$ captures a positive fraction of $\lvert Q \rvert^2$.  The only
analytic input is a Poincar\'e inequality quantifying how well $Q$ is
approximated by its local averages.

\begin{theorem}[High-frequency band estimate]\label{thm:high-freq}
\par\noindent\texttt{FockSPR.high\_freq\_band\_estimate}\par\smallskip
Let $N \ge 1$ and $L \ge 1$ satisfy $1343\, L^2 \le N^2$.
For any $b_0, \dots, b_{L-1} \in \mathbb{C}$, the polynomial
$P(t) = \sum_{m=0}^{L-1} b_m\, \mathrm{e}^{\mathrm{i}(N+m)t}$ satisfies
\[
\lVert P \rVert_{L^2}^2
\;\le\;
32\, \lVert \rho \circ P \rVert_{L^2}^2.
\]
\end{theorem}

\begin{proof}
Write $P(t) = \mathrm{e}^{\mathrm{i} Nt}\,Q(t)$ with
$Q(t) = \sum_{m=0}^{L-1} b_m\,\mathrm{e}^{\mathrm{i} mt}$, so $\lVert P \rVert = \lVert Q \rVert$.
Partition the circle into $N$ equal arcs
$J_k = [2\pi k/N,\; 2\pi(k+1)/N)$, and let
$q_k = N\!\int_{J_k}\! Q\,\mathrm{d}\sigma$ be the average of $Q$ on~$J_k$.

\medskip
\noindent
\emph{Step 1: Poincar\'e approximation.}\;
On each interval $J_k$ of length $h = 2\pi/N$, the Poincar\'e
inequality gives
$\int_{J_k}\lvert Q - q_k \rvert^2\,\mathrm{d}\sigma \le h^2\int_{J_k}\lvert Q' \rvert^2\,
\mathrm{d}\sigma$.\footnote{We use the inequality
$\int_0^h \lvert f - \bar f \rvert^2\,\mathrm{d} x
\le h^2\int_0^h\lvert f' \rvert^2\,\mathrm{d} x$
with the non-sharp constant~$1$; here $\bar f$ denotes the mean of~$f$
on~$[0,h]$.  The sharp Wirtinger
constant is $1/\pi^2$, which would improve the threshold from $1343$ to~$136$.
We use this weaker version because it has a short proof from Jensen's
inequality and Cauchy--Schwarz, and we formalised this proof in Lean; the
sharp Wirtinger inequality is not currently available in Mathlib.}
Summing over $k$ and using Bernstein's inequality $\lVert Q' \rVert^2 \le L^2\lVert Q \rVert^2$
(since every frequency of~$Q$ has modulus at most~$L{-}1 < L$):
\begin{equation}\label{eq:delta}
\delta \;:=\; \sum_{k} \int_{J_k}\lvert Q - q_k \rvert^2\,\mathrm{d}\sigma
\;\le\; \frac{4\pi^2 L^2}{N^2}\,\lVert Q \rVert^2.
\end{equation}

\noindent
\emph{Step 2: Parseval partition.}\;
The mean-variance decomposition on each $J_k$ gives
\begin{equation}\label{eq:parseval-partition}
\lVert Q \rVert^2
= \frac{1}{N}\sum_{k}\lvert q_k \rvert^2 + \delta.
\end{equation}

\noindent
\emph{Step 3: Frozen rotation.}\;
As $t$ traverses $J_k$, the argument $Nt$ traverses one full period.
Substituting $s = Nt$ and applying Proposition~\ref{prop:rot-avg} with
$r = \lvert q_k \rvert$:
\[
\int_{J_k}\rho\bigl(\mathrm{e}^{\mathrm{i} Nt} q_k\bigr)^2\,\mathrm{d}\sigma
= \frac{1}{N}\int_{S^1}\rho\bigl(\lvert q_k \rvert\mathrm{e}^{\mathrm{i} s}\bigr)^2\,\mathrm{d}\sigma(s)
\;\ge\; \frac{\lvert q_k \rvert^2}{8N}.
\]

\noindent
\emph{Step 4: Assembly.}\;
On $J_k$, the Lipschitz bound~\eqref{eq:rho-lip} gives
$\rho(P) = \rho(\mathrm{e}^{\mathrm{i} Nt}Q)
\ge \rho(\mathrm{e}^{\mathrm{i} Nt}q_k) - \lvert Q - q_k \rvert$.
By Lemma~\ref{lem:safe-square},
$\rho(P)^2 \ge \tfrac{1}{2}\rho(\mathrm{e}^{\mathrm{i} Nt}q_k)^2 - \lvert Q - q_k \rvert^2$.
Integrating over $J_k$ and using Step~3:
\[
\int_{J_k}\rho(P)^2\,\mathrm{d}\sigma
\;\ge\; \frac{\lvert q_k \rvert^2}{16N} - \int_{J_k}\lvert Q-q_k \rvert^2\,\mathrm{d}\sigma.
\]
Summing over $k$ and applying~\eqref{eq:parseval-partition}:
\[
\lVert \rho \circ P \rVert^2
\;\ge\; \frac{1}{16}\Bigl(\lVert Q \rVert^2 - \delta\Bigr) - \delta
= \frac{\lVert Q \rVert^2}{16} - \frac{17\delta}{16}.
\]
Substituting the bound~\eqref{eq:delta}:
\[
\lVert \rho \circ P \rVert^2
\;\ge\;
\frac{1}{16}\Bigl(1 - \frac{68\pi^2 L^2}{N^2}\Bigr)\lVert Q \rVert^2.
\]
The condition $1343\,L^2 \le N^2$ gives $68\pi^2 L^2/N^2 \le 1/2$
(since $136\pi^2 < 1343$), and therefore
$\lVert \rho \circ P \rVert^2 \ge \lVert P \rVert^2/32$.
\end{proof}
 \section{Frequency localisation}\label{sec:blocks}

The circle estimates of Section~\ref{sec:circle} require the trigonometric
polynomial to have a controlled number of frequencies, yet $U$ may have
arbitrarily many.  The Gaussian weight resolves this: the radial factor
$r \mapsto r^{2n+1}\mathrm{e}^{-r^2}$ peaks sharply near $r = \sqrt{n + 1/2}$,
so grouping frequencies by peak location decomposes $U$ into ``blocks''
approximately supported on disjoint annuli.

\subsection{Blocks and their properties}\label{sec:block-def}

For $\ell \ge 1$, define the \emph{frequency block}
\[
I_\ell = \bigl\{n \in \mathbb{N} : \ell^2 \le n < (\ell+1)^2\bigr\},
\]
which has $\lvert I_\ell \rvert = 2\ell+1$ elements, and the corresponding
\emph{block polynomial}
$U_\ell(z) = \sum_{n \in I_\ell,\, n \le D} a_n z^n$.
\par\noindent\texttt{FockSPR.freqBlock, FockSPR.blockPoly}\par\smallskip
The blocks $I_1, I_2, \dots$ are pairwise disjoint and cover
$\{1, \dots, D\}$ up to index $\Lambda = \lfloor\sqrt{D}\rfloor$.  Since the
Fock norm is a sum over individual frequencies, it decomposes accordingly:
\begin{equation}\label{eq:fock-decomp}
\lVert U \rVert_\mathcal{F}^2 = \sum_{\ell=1}^{\Lambda} \lVert U_\ell \rVert_\mathcal{F}^2.
\end{equation}
\par\noindent\texttt{FockSPR.fockNorm\_decomposes}\par\smallskip
Similarly, on each circle of radius~$r$, blocks with distinct frequency
supports are orthogonal:
$\int_{S^1} U_{\ell_1}(r\mathrm{e}^{\mathrm{i} t})\,
\overline{U_{\ell_2}(r\mathrm{e}^{\mathrm{i} t})}\,\mathrm{d}\sigma(t) = 0$
for $\ell_1 \ne \ell_2$.

The key geometric fact is that \emph{each block has its peaks in the
corresponding annulus}:

\begin{lemma}[Peak localisation]\label{lem:peak}
\par\noindent\texttt{FockSPR.monomial\_peak\_localization}\par\smallskip
For $n \in I_\ell$ with $\ell \ge 1$, the saddle point
$r_n = \sqrt{n + 1/2}$ satisfies $\ell < r_n < \ell + 1$.
\end{lemma}

\begin{proof}
From $n \ge \ell^2$: $r_n > \ell$.  From $n < (\ell+1)^2$:
$r_n < \ell + 1$.
\end{proof}

\subsection{Local and remainder polynomials}\label{sec:local-remainder}

Fix a window width $M \ge 1$.  For each annulus index $j \ge 0$, define the
\emph{local polynomial}
\[
V_j = \sum_{\substack{1 \le \ell \le \Lambda \\ \lvert \ell - j \rvert \le M}} U_\ell
\qquad\text{and the \emph{remainder}}\qquad
R_j = U - V_j.
\]
\par\noindent\texttt{FockSPR.localPoly}\par\smallskip
The local polynomial gathers the blocks whose peaks fall within $M$ steps of
annulus~$j$; the remainder gathers the rest.

For $j \ge M + 1$, the frequency support of $V_j$ lies in
$[(j{-}M)^2,\; (j{+}M{+}1)^2)$, and the number of frequencies is
$L_j = (2M{+}1)(2j{+}1)$.

\subsection{Saddle-point estimates for the leakage}\label{sec:saddle}

The energy that a block $U_\ell$ contributes to a distant annulus decays
like $\mathrm{e}^{-(\lvert j-\ell \rvert-1)^2}$.  The mechanism is simple: the function
$\varphi_n(r) = (2n{+}1)\log r - r^2$ (so that
$r^{2n+1}\mathrm{e}^{-r^2} = \mathrm{e}^{\varphi_n(r)}$) has a unique maximum at
$r_n = \sqrt{n+1/2}$ and concavity $\varphi_n''(r) \le -2$, giving the
Gaussian upper bound
\begin{equation}\label{eq:phi-quad}
\mathrm{e}^{\varphi_n(r)}
\;\le\;
\mathrm{e}^{\varphi_n(r_n)}\, \mathrm{e}^{-(r - r_n)^2}.
\end{equation}
\par\noindent\texttt{FockSPR.phiFunc\_quad\_bound}\par\smallskip
The value at the saddle point is controlled by the factorial:

\begin{lemma}\label{lem:exp-phi}
\par\noindent\texttt{FockSPR.exp\_phi\_le\_factorial}\par\smallskip
For $n \ge 1$,\;
$\mathrm{e}^{\varphi_n(r_n)} \le \mathrm{e}^{1/4}\, n!/2$.
\end{lemma}

\begin{proof}
The Gamma integral gives
$n!/2 = \int_0^\infty \mathrm{e}^{\varphi_n(r)}\,\mathrm{d} r$.
Since $r_n \ge \sqrt{3/2} > 1$, the interval $[r_n - 1/2,\, r_n + 1/2]$
lies in $(0,\infty)$, and~\eqref{eq:phi-quad} shows that
$\varphi_n \ge \varphi_n(r_n) - 1/4$ on this interval.  Hence
$n!/2 \ge \mathrm{e}^{\varphi_n(r_n) - 1/4}$.
\end{proof}

Combining~\eqref{eq:phi-quad} and Lemma~\ref{lem:exp-phi}: for $r \in
[j, j{+}1]$ and $n \in I_\ell$, the monomial weight satisfies
\begin{equation}\label{eq:monomial-bound}
r^{2n+1}\mathrm{e}^{-r^2}
\;\le\;
\frac{\mathrm{e}^{1/4}\, n!}{2}\,
\mathrm{e}^{-\operatorname{dist}(r_n,\,[j,j{+}1])^2},
\end{equation}
where $\operatorname{dist}(r_n, [j,j{+}1]) = \max(j - r_n,\, r_n - (j{+}1),\, 0)$.
By Lemma~\ref{lem:peak}, $r_n \in (\ell, \ell{+}1)$, so this distance is at
least $(\lvert j - \ell \rvert - 1)_+$.

\subsection{Leakage bounds}

\begin{proposition}[Single-block leakage]\label{prop:single-leak}
\par\noindent\texttt{FockSPR.block\_annulus\_leakage}\par\smallskip
For $\ell \ge 1$ and $j \ge 0$,
\[
\frac{1}{\pi}\int_{\{j \le \lvert z \rvert < j+1\}}
\lvert U_\ell(z) \rvert^2\, \mathrm{e}^{-\lvert z \rvert^2}\,\mathrm{d} m(z)
\;\le\;
\mathrm{e}^{1/4}\,
\mathrm{e}^{-(\lvert j-\ell \rvert-1)_+^2}\,
\lVert U_\ell \rVert_\mathcal{F}^2.
\]
\end{proposition}

\begin{proof}
Passing to polar coordinates and using the orthogonality of distinct
monomials on each circle, the left side equals
$2\sum_{n \in I_\ell}\lvert a_n \rvert^2 \int_j^{j+1} r^{2n+1}\mathrm{e}^{-r^2}\mathrm{d} r$.
Bounding each radial integral by~\eqref{eq:monomial-bound} and summing:
$\le \mathrm{e}^{1/4}\mathrm{e}^{-(\lvert j-\ell \rvert-1)_+^2}\sum_{n \in I_\ell}\lvert a_n \rvert^2 n!$.
\end{proof}

Summing over all distant blocks and swapping the order of summation
yields:

\begin{proposition}[Total leakage]\label{prop:total-leak}
\par\noindent\texttt{FockSPR.total\_leakage\_bound}\par\smallskip
For $M \ge 2$, define $\eta_M = 2\mathrm{e}^{1/4}\sum_{m \ge M}\mathrm{e}^{-m^2}$.  Then
\[
\sum_{j \ge 0}\;
\frac{1}{\pi}\int_{\{j \le \lvert z \rvert < j+1\}}
\lvert R_j(z) \rvert^2\, \mathrm{e}^{-\lvert z \rvert^2}\,\mathrm{d} m(z)
\;\le\;
\eta_M \cdot \lVert U \rVert_\mathcal{F}^2.
\]
\end{proposition}

\begin{proof}
By circle orthogonality on each annulus,
$\int \lvert R_j \rvert^2 = \sum_{\lvert \ell-j \rvert>M}\int\lvert U_\ell \rvert^2$.
Apply Proposition~\ref{prop:single-leak}, swap $\sum_j$ and $\sum_\ell$
(all terms are nonneg\-ative), and note that for $\lvert j-\ell \rvert > M \ge 2$
the decay factor is $\mathrm{e}^{-(\lvert j-\ell \rvert-1)^2}$.  The inner sum over
$j$ contributes at most $2\sum_{m \ge M}\mathrm{e}^{-m^2}$; multiplying by
$\mathrm{e}^{1/4}$ and using~\eqref{eq:fock-decomp} gives the result.
\end{proof}

For $M = 5$: the geometric series bound
$\sum_{m \ge 5}\mathrm{e}^{-m^2} \le \mathrm{e}^{-25}/(1 - \mathrm{e}^{-11}) < 1.39 \times
10^{-11}$ gives:

\begin{corollary}\label{cor:eta5}
\par\noindent\texttt{FockSPR.eta\_5\_bound}\par\smallskip
$\eta_5 < 4 \times 10^{-11}$.
\end{corollary}

\subsection{The annulus estimate}\label{sec:annulus}

We can now state the uniform estimate that holds on every annulus.  The
function $\zeta \mapsto V_j(r\zeta)$ is a trigonometric polynomial on $S^1$
with all frequencies in $\mathbb{Z}_{>0}$, and we apply either
Theorem~\ref{thm:local-circle} or Theorem~\ref{thm:high-freq} depending
on~$j$.

\begin{proposition}[Annulus local estimate]\label{prop:annulus}
\par\noindent\texttt{FockSPR.annulus\_local\_estimate}\par\smallskip
With $M = 5$: for every $j \ge 0$ and every $r \in [j, j{+}1]$,
\[
\int_{S^1}\lvert V_j(r\zeta) \rvert^2\,\mathrm{d}\sigma(\zeta)
\;\le\;
1620^2 \cdot
\int_{S^1}\rho\bigl(V_j(r\zeta)\bigr)^2\,\mathrm{d}\sigma(\zeta).
\]
\end{proposition}

\begin{proof}
When $r = 0$ both sides vanish (all monomials have positive degree).
For $r > 0$, the function $\zeta \mapsto V_j(r\zeta)$ is a trigonometric
polynomial with positive frequencies and coefficients $a_n r^n$.

\smallskip\noindent
\emph{Case $j \le 817$.}\;
With $M = 5$, the frequency count satisfies $L_j \le 11(2 \cdot 817 + 1)
< 18{,}000$.  By Theorem~\ref{thm:local-circle},
$\lVert V_j(r\,\cdot\,) \rVert^2 \le 144\, L_j \cdot
\lVert \rho \circ V_j(r\,\cdot\,) \rVert^2$, and
$144 \cdot 18{,}000 < 1620^2$.

\smallskip\noindent
\emph{Case $j \ge 818$.}\;
The lowest frequency is $N_j = (j{-}5)^2$ and the bandwidth is
$L_j = 11(2j{+}1)$.  The ratio $L_j^2/N_j^2$ is decreasing in~$j$
(for $j \ge 6$), so it suffices to verify the condition
$1343\,L_j^2 \le N_j^2$ of Theorem~\ref{thm:high-freq} at $j = 818$.
A direct computation confirms $1343 \cdot (11 \cdot 1637)^2 < 813^4$.
The high-frequency estimate then gives
$\lVert V_j(r\,\cdot\,) \rVert^2 \le 32\,\lVert \rho \circ V_j(r\,\cdot\,) \rVert^2
\le 1620^2\,\lVert \rho \circ V_j(r\,\cdot\,) \rVert^2$.
\end{proof}
 \section{Proof of the orthogonal coercivity estimate}\label{sec:assembly}

With the circle estimates (Section~\ref{sec:circle}), the leakage bounds
(Section~\ref{sec:blocks}), and the annulus estimate
(Proposition~\ref{prop:annulus}) in hand, the proof of
Theorem~\ref{thm:orthog} is a short calculation.

\begin{proof}[Proof of Theorem~\ref{thm:orthog}]
Fix $M = 5$ throughout, and write $B = 1620$ for the annulus constant from
Proposition~\ref{prop:annulus}.  Let $V_j$ and $R_j$ be the local and
remainder polynomials from Section~\ref{sec:local-remainder}.

\medskip
\noindent
\emph{Step 1: From $V_j$ to $\rho(U)$ on each circle.}\;
For $r \in [j, j{+}1]$, Proposition~\ref{prop:annulus} gives
\begin{equation}\label{eq:annulus-applied}
\int_{S^1}\lvert V_j(r\zeta) \rvert^2\,\mathrm{d}\sigma
\;\le\;
B^2\int_{S^1}\rho\bigl(V_j(r\zeta)\bigr)^2\,\mathrm{d}\sigma.
\end{equation}
Since $\rho$ is $1$-Lipschitz and $U = V_j + R_j$,
\[
\rho\bigl(V_j(z)\bigr)
\le \rho\bigl(U(z)\bigr) + \lvert R_j(z) \rvert.
\]
Squaring ($(a+b)^2 \le 2a^2 + 2b^2$) and substituting
into~\eqref{eq:annulus-applied}:
\begin{equation}\label{eq:Vj-to-U}
\int_{S^1}\lvert V_j(r\zeta) \rvert^2\,\mathrm{d}\sigma
\;\le\;
2B^2\!\int_{S^1}\rho(U(r\zeta))^2\,\mathrm{d}\sigma
\;+\; 2B^2\!\int_{S^1}\lvert R_j(r\zeta) \rvert^2\,\mathrm{d}\sigma.
\end{equation}

\medskip
\noindent
\emph{Step 2: From $U$ to $V_j$ on each circle.}\;
Applying $\lvert U \rvert^2 \le 2\lvert V_j \rvert^2 + 2\lvert R_j \rvert^2$
and~\eqref{eq:Vj-to-U}:
\begin{equation}\label{eq:circle-ineq}
\int_{S^1}\lvert U(r\zeta) \rvert^2\,\mathrm{d}\sigma
\;\le\;
4B^2\!\int_{S^1}\rho(U)^2\,\mathrm{d}\sigma
\;+\; (4B^2 + 2)\!\int_{S^1}\lvert R_j \rvert^2\,\mathrm{d}\sigma.
\end{equation}

\medskip
\noindent
\emph{Step 3: Integrate over annuli.}\;
Multiply~\eqref{eq:circle-ineq} by $2r\mathrm{e}^{-r^2}$, integrate $r$ over
$[j, j{+}1]$, and use the polar decomposition~\eqref{eq:polar}.  The
$\rho(U)$ term on the right does not depend on~$j$ (the annuli tile the
plane), while the $R_j$ term does.  Summing over $j \ge 0$:
\[
\lVert U \rVert_\mathcal{F}^2
\;\le\;
4B^2\,\lVert \rho(U) \rVert_\mathcal{F}^2
\;+\;
(4B^2 + 2)\sum_{j \ge 0}\;
\frac{1}{\pi}\int_{\{j \le \lvert z \rvert < j+1\}}\!\!
\lvert R_j \rvert^2\mathrm{e}^{-\lvert z \rvert^2}\mathrm{d} m.
\]
By Proposition~\ref{prop:total-leak} with $M = 5$, the remainder sum is at
most $\eta_5\,\lVert U \rVert_\mathcal{F}^2$.  Hence
\begin{equation}\label{eq:absorb}
\lVert U \rVert_\mathcal{F}^2
\;\le\;
4B^2\,\lVert \rho(U) \rVert_\mathcal{F}^2
\;+\;
(4B^2 + 2)\,\eta_5\,\lVert U \rVert_\mathcal{F}^2.
\end{equation}

\medskip
\noindent
\emph{Step 4: Absorption.}\;
Since $4B^2 + 2 < 10^7$ and $\eta_5 < 4 \times 10^{-11}$, the coefficient
of $\lVert U \rVert_\mathcal{F}^2$ on the right of~\eqref{eq:absorb} satisfies
$(4B^2 + 2)\,\eta_5 < 4 \times 10^{-4} < \tfrac{1}{2}$.
Moving this term to the left and using $1 - (4B^2+2)\eta_5 > 1/2$:
\[
\lVert U \rVert_\mathcal{F}^2
\;\le\; 8B^2\,\lVert \rho(U) \rVert_\mathcal{F}^2
\;=\; 8 \cdot 1620^2\,\lVert \rho(U) \rVert_\mathcal{F}^2
\;\le\; 4600^2\,\lVert \rho(U) \rVert_\mathcal{F}^2.\qedhere
\]
\end{proof}

\begin{remark}[On the constant]
The bottleneck is the crossover at $j^* = 818$ between the two circle
estimates.  Below this threshold, the local polynomial has up to
$L \approx 18{,}000$ frequencies, producing the large annulus constant
$B = 1620$; the leakage and the squaring step ($8B^2 \le C^2$) contribute
only modest additional factors.

The single most effective improvement would be to replace the weak
Poincar\'e constant in Theorem~\ref{thm:high-freq} by the sharp Wirtinger
constant $1/\pi^2$.  This would lower the high-frequency threshold from
$1343\,L^2 \le N^2$ to $136\,L^2 \le N^2$, moving the crossover to
$j^* \approx 90$ and reducing $C$ to approximately~$3700$.
\end{remark}

\begin{remark}[Formalisation]
The entire proof---including the orthogonal reduction
(Theorem~\ref{thm:reduction}), the orthogonal coercivity estimate
(Theorem~\ref{thm:orthog}), and their combination into the local stability
theorem (Theorem~\ref{thm:local})---has been formalised in Lean~4 using the
Mathlib library~\cite{mathlib}.  The formalisation closely mirrors the paper
structure, and the theorem names recorded in blue monospace
throughout can be searched directly in the accompanying repository.

The only result \emph{not} taken from Mathlib is the weak Poincar\'e
inequality on intervals used in the proof of
Theorem~\ref{thm:high-freq} (see the footnote there), which we proved
from scratch using Jensen's inequality and Cauchy--Schwarz.
All other analytic ingredients---polar coordinates, the Gamma integral,
Fourier orthonormality, Parseval's identity---are imported from Mathlib.
\end{remark}
 
\begin{thebibliography}{10}

\bibitem{alaifari2017continuous}
R.~Alaifari and P.~Grohs.
\newblock Phase retrieval in the general setting of continuous frames for
  {B}anach spaces.
\newblock \emph{SIAM J.\ Math.\ Anal.}, 49(3):1895--1911, 2017.

\bibitem{alaifari2021gabor}
R.~Alaifari and P.~Grohs.
\newblock Gabor phase retrieval is severely ill-posed.
\newblock \emph{Appl.\ Comput.\ Harmon.\ Anal.}, 50:401--419, 2021.

\bibitem{bargmann1961hilbert}
V.~Bargmann.
\newblock On a {H}ilbert space of analytic functions and an associated integral
  transform.
\newblock \emph{Comm.\ Pure Appl.\ Math.}, 14:187--214, 1961.

\bibitem{calderbank2022stable}
R.~Calderbank, I.~Daubechies, D.~Freeman, and N.~Freeman.
\newblock Stable phase retrieval for infinite dimensional subspaces of
  $L_2(\mathbb{R})$.
\newblock arXiv:2203.03135, 2022.

\bibitem{christ2022examples}
M.~Christ, B.~Pineau, and M.~A.~Taylor.
\newblock Examples of {H}\"older-stable phase retrieval.
\newblock \emph{Math.\ Res.\ Lett.}, 31(5):1339--1352, 2024.

\bibitem{grohs2020survey}
P.~Grohs, S.~Koppensteiner, and M.~Rathmair.
\newblock Phase retrieval: uniqueness and stability.
\newblock \emph{SIAM Rev.}, 62(2):301--350, 2020.

\bibitem{FOTP}
D.~Freeman, T.~Oikhberg, B.~Pineau, and M.~A.~Taylor.
\newblock Stable phase retrieval in function spaces.
\newblock \emph{Math.\ Ann.}, 390(1):1--43, 2024.

\bibitem{mathlib}
The Mathlib Community.
\newblock \emph{Mathlib: the Lean~4 mathematics library}.
\newblock \url{https://github.com/leanprover-community/mathlib4}, 2024.

\end{thebibliography}

\end{document}
